% This document was generated by an LLM.

\documentclass{essay}

\newcommand{\F}{\mathbf{F}}
\newcommand{\rr}{\mathbf{r}}
\newcommand{\nn}{\mathbf{n}}
\newcommand{\TT}{\mathbf{T}}

\newcommand{\grad}{\operatorname{grad}}
\newcommand{\dive}{\operatorname{div}}
\newcommand{\curl}{\operatorname{curl}}

\title{Types in Multivariable Calculus}
\subtitle{Reading the Notation of Line, Surface, and Vector Integrals}
\author{}
\date{}

\begin{document}

\maketitle

\begin{abstract}
\noindent
The notation of vector calculus is dense because it hides the \emph{types} of its
objects. A single symbol like $\dd\mathbf{S}$ or $\nabla\cdot\mathbf{F}$ silently
declares whether something is a number or a vector, and whether an operation
consumes a field and returns a field or a scalar. This note treats every object as
if it had a function signature: domain in, codomain out. We unpack each
differential element ($\dd s$, $\dd\mathbf{r}$, $\dd S$, $\dd\mathbf{S}$, $\dd V$),
each operator ($\nabla$, $\nabla\cdot$, $\nabla\times$), and each integral, and we
show that every one of them collapses, after a parametrization, into an ordinary
single-variable Riemann integral. The recurring theme: \emph{a parametrization is a
type cast} that turns a geometric integral into the $\int_a^b(\cdots)\,\dd t$ you
already know.
\end{abstract}

\tableofcontents

\section{The typing problem}

A function signature tells you, before you read a single line of the body, what goes
in and what comes out: \texttt{length : string -> int}. You know immediately that
feeding it an \texttt{int} is an error, and that the result is a number you can add
to other numbers.

Vector-calculus notation has the opposite ergonomics. The expression
\[
\iint_S \F\cdot \dd\mathbf{S}
\]
declares its types only implicitly. To read it you must already know that $S$ is a
\emph{surface}, that $\F$ is a \emph{vector field}, that $\dd\mathbf{S}$ is an
\emph{infinitesimal vector} (not a scalar), that ``$\cdot$'' is the dot product
(which forces the result of the integrand to be a \emph{scalar}), and therefore that
the whole expression evaluates to a single \emph{number}. None of that is written
down; you reconstruct it in your head.

This document makes the reconstruction explicit. For every piece of notation we
state, in the style of a function signature, what type each symbol carries and what
type the combination produces.

\subsection{The type vocabulary}

We will work in $\R^2$ or $\R^3$. The objects fall into a small number of types.

\begin{center}
\renewcommand{\arraystretch}{1.35}
\begin{tabular}{@{}l l l@{}}
\toprule
\textbf{Type} & \textbf{Notation} & \textbf{What it is} \\
\midrule
point          & $\mathbf{p}=(x,y,z)\in\R^3$        & a location in space \\
scalar         & $c\in\R$                           & a single number \\
vector         & $\mathbf{v}\in\R^3$                & a magnitude-and-direction arrow \\
scalar field   & $f:\R^3\to\R$                      & a number at every point \\
vector field   & $\F:\R^3\to\R^3$                   & an arrow at every point \\
parametrized curve   & $\rr:[a,b]\to\R^3$          & a $1$-parameter path \\
parametrized surface & $\rr:D\subset\R^2\to\R^3$   & a $2$-parameter sheet \\
\bottomrule
\end{tabular}
\end{center}

Two distinctions do almost all the work:

\begin{itemize}[leftmargin=1.4em]
\item \textbf{scalar vs.\ vector.} A scalar field assigns a \emph{number}
(temperature, pressure, potential). A vector field assigns an \emph{arrow}
(velocity, force, flow). Confusing the two is the single most common reading error.
\item \textbf{point vs.\ vector.} A point is \emph{where} you are; a vector is a
\emph{displacement} or arrow. They live in the same $\R^3$ and are written the same
way, which is exactly why the notation is slippery. A field eats a point and returns
either a scalar or a vector.
\end{itemize}

\section{Differential elements, built from single-variable calculus}

Every integral below is glued together from an \emph{infinitesimal element}: $\dd s$,
$\dd\mathbf{r}$, $\dd S$, $\dd\mathbf{S}$, or $\dd V$. Each one has a type, and each
one is derived by the same move you already use in one dimension, namely
$\dd u = u'(t)\,\dd t$.

\subsection{The scalar arc-length element $\dd s$}

\textbf{Type: scalar.} $\dd s$ is an infinitesimal \emph{length}, a non-negative
number.

Start from the Pythagorean theorem applied to an infinitesimal right triangle whose
legs are $\dd x$ and $\dd y$:
\[
\dd s = \sqrt{\dd x^2 + \dd y^2}
\qquad\bigl(\text{in }\R^3:\ \dd s=\sqrt{\dd x^2+\dd y^2+\dd z^2}\bigr).
\]
This is just $a^2+b^2=c^2$ for the hypotenuse of a tiny triangle. To make it
computable, introduce a parameter $t$ with $\rr(t)=(x(t),y(t))$. Then each
differential is an ordinary single-variable differential,
\[
\dd x = x'(t)\,\dd t, \qquad \dd y = y'(t)\,\dd t,
\]
and substituting,
\[
\dd s
= \sqrt{x'(t)^2\,\dd t^2 + y'(t)^2\,\dd t^2}
= \sqrt{x'(t)^2 + y'(t)^2}\;\dd t
= \lvert \rr'(t)\rvert \,\dd t .
\]
The factor $\lvert\rr'(t)\rvert$ is the \emph{speed}: a scalar. This is the whole
trick in miniature. The geometric object $\dd s$ becomes ``speed times $\dd t$'',
which is exactly the kind of thing you integrate in first-year calculus.

\subsection{The vector line element $\dd\mathbf{r}$}

\textbf{Type: vector.} $\dd\mathbf{r}$ is an infinitesimal \emph{displacement
arrow}, pointing along the curve in the direction of increasing parameter.

Where $\dd s$ keeps only the length, $\dd\mathbf{r}$ keeps the components:
\[
\dd\mathbf{r} = (\dd x,\ \dd y,\ \dd z)
= \bigl(x'(t),\,y'(t),\,z'(t)\bigr)\,\dd t
= \rr'(t)\,\dd t .
\]
So $\dd\mathbf{r}$ is the velocity vector $\rr'(t)$ scaled by $\dd t$. Its length is
$\lvert\dd\mathbf{r}\rvert = \lvert\rr'(t)\rvert\,\dd t = \dd s$, which gives the
bridge between the two:
\[
\boxed{\ \dd\mathbf{r} = \TT\,\dd s\ }
\qquad\text{where}\qquad
\TT = \frac{\rr'(t)}{\lvert\rr'(t)\rvert}\ \ \text{is the unit tangent (a vector).}
\]
Read the boxed identity as a type statement: \emph{vector} $\dd\mathbf{r}$ equals
\emph{unit vector} $\TT$ times \emph{scalar} $\dd s$. The unit vector carries the
direction; the scalar carries the length.

\subsection{The scalar area element $\dd S$}

\textbf{Type: scalar.} $\dd S$ is an infinitesimal \emph{area}, a non-negative
number.

For a curve we measured the length of one tangent vector. For a surface we measure
the \emph{area of the parallelogram} spanned by two tangent vectors. Parametrize the
surface by two variables, $\rr(u,v)$. Moving $u$ a little traces the tangent vector
$\rr_u\,\dd u$; moving $v$ a little traces $\rr_v\,\dd v$. The area of the
parallelogram they span is the magnitude of their cross product:
\[
\dd S = \bigl\lvert \rr_u \times \rr_v \bigr\rvert \,\dd u\,\dd v .
\]
Here $\rr_u=\partial\rr/\partial u$ and $\rr_v=\partial\rr/\partial v$ are vectors,
their cross product is a vector, and its magnitude is a scalar. Compare directly:
\[
\underbrace{\dd s = \lvert\rr'(t)\rvert\,\dd t}_{\text{length of a tangent vector}}
\qquad\longleftrightarrow\qquad
\underbrace{\dd S = \lvert\rr_u\times\rr_v\rvert\,\dd u\,\dd v}_{\text{area of a tangent parallelogram}} .
\]
One dimension uses the magnitude of a derivative; two dimensions use the magnitude
of a cross product of partial derivatives. Same idea, one rung up.

\subsection{The vector area element $\dd\mathbf{S}$}

\textbf{Type: vector.} $\dd\mathbf{S}$ is an infinitesimal area carrying a
\emph{direction}: it points perpendicular to the surface, with magnitude equal to
$\dd S$.

This is the surface analogue of $\dd\mathbf{r}=\TT\,\dd s$. The cross product
$\rr_u\times\rr_v$ is already perpendicular to the surface, so we keep it whole
instead of taking its magnitude:
\[
\dd\mathbf{S} = (\rr_u \times \rr_v)\,\dd u\,\dd v = \nn\,\dd S,
\qquad
\nn = \frac{\rr_u\times\rr_v}{\lvert\rr_u\times\rr_v\rvert}\ \ \text{(unit normal).}
\]
Again read it as types: \emph{vector} $\dd\mathbf{S}$ equals \emph{unit vector}
$\nn$ times \emph{scalar} $\dd S$. The normal $\nn$ is what makes ``flux through a
surface'' meaningful: it records which way you are passing through.

\begin{center}
\renewcommand{\arraystretch}{1.35}
\begin{tabular}{@{}l l l l@{}}
\toprule
\textbf{Geometry} & \textbf{Scalar element} & \textbf{Vector element} & \textbf{Link} \\
\midrule
curve   & $\dd s = \lvert\rr'\rvert\,\dd t$
        & $\dd\mathbf{r}=\rr'\,\dd t$
        & $\dd\mathbf{r}=\TT\,\dd s$ \\
surface & $\dd S=\lvert\rr_u\times\rr_v\rvert\,\dd u\,\dd v$
        & $\dd\mathbf{S}=(\rr_u\times\rr_v)\,\dd u\,\dd v$
        & $\dd\mathbf{S}=\nn\,\dd S$ \\
\bottomrule
\end{tabular}
\end{center}

\subsection{The volume element $\dd V$}

\textbf{Type: scalar.} $\dd V$ is an infinitesimal \emph{volume}. In Cartesian
coordinates $\dd V = \dd x\,\dd y\,\dd z$. Under a change of variables
$\rr(u,v,w)$ it picks up the absolute value of the Jacobian determinant,
\[
\dd V = \bigl\lvert \det[\,\rr_u\ \rr_v\ \rr_w\,] \bigr\rvert \,\dd u\,\dd v\,\dd w,
\]
which is the three-dimensional version of the parallelogram-area idea: the
determinant of the three tangent vectors is the (signed) volume of the
parallelepiped they span.

\section{The three operators as typed functions}

The symbols $\nabla$, $\nabla\cdot$, and $\nabla\times$ look like three unrelated
notations. They are better understood as one object, $\nabla$, combined with the
three products a vector can participate in. Write $\nabla$ as a ``vector of
partial-derivative operators'',
\[
\nabla = \Bigl(\frac{\partial}{\partial x},\ \frac{\partial}{\partial y},\
\frac{\partial}{\partial z}\Bigr).
\]
This is a type pun, but a productive one: $\nabla$ behaves like a vector whose
components happen to be operators.

\begin{center}
\renewcommand{\arraystretch}{1.5}
\begin{tabular}{@{}l l l l@{}}
\toprule
\textbf{Operator} & \textbf{Product pattern} & \textbf{Signature} & \textbf{Result type} \\
\midrule
$\grad f = \nabla f$           & scalar $\times$ vector & $(\R^3\!\to\!\R)\to(\R^3\!\to\!\R^3)$ & vector field \\
$\dive \F = \nabla\cdot\F$     & vector $\cdot$ vector  & $(\R^3\!\to\!\R^3)\to(\R^3\!\to\!\R)$ & scalar field \\
$\curl \F = \nabla\times\F$    & vector $\times$ vector & $(\R^3\!\to\!\R^3)\to(\R^3\!\to\!\R^3)$ & vector field \\
\bottomrule
\end{tabular}
\end{center}

The result type is forced by the product, exactly as in ordinary vector algebra:
a scalar times a vector is a vector; a dot product of two vectors is a scalar; a
cross product of two vectors is a vector. If you remember the output types of the
three products, you remember the output types of grad, div, and curl for free.

\subsection{Component formulas}

Let $f$ be a scalar field and $\F=(P,Q,R)$ a vector field, with $P,Q,R:\R^3\to\R$.

\textbf{Gradient} (scalar field $\to$ vector field). Each component is a partial
derivative; the field of steepest ascent:
\[
\nabla f = \Bigl(\frac{\partial f}{\partial x},\ \frac{\partial f}{\partial y},\
\frac{\partial f}{\partial z}\Bigr).
\]

\textbf{Divergence} (vector field $\to$ scalar field). The dot product
$\nabla\cdot\F$ sums the matching partials; it measures net outflow per unit volume:
\[
\nabla\cdot\F = \frac{\partial P}{\partial x} + \frac{\partial Q}{\partial y}
+ \frac{\partial R}{\partial z}.
\]

\textbf{Curl} (vector field $\to$ vector field). The cross product
$\nabla\times\F$ is the formal determinant
\[
\nabla\times\F =
\begin{vmatrix}
\mathbf{i} & \mathbf{j} & \mathbf{k} \\[2pt]
\dfrac{\partial}{\partial x} & \dfrac{\partial}{\partial y} & \dfrac{\partial}{\partial z} \\[6pt]
P & Q & R
\end{vmatrix}
=\Bigl(R_y - Q_z,\ \ P_z - R_x,\ \ Q_x - P_y\Bigr),
\]
where subscripts denote partial derivatives. It measures local rotation.

\textbf{A note on dimension.} Divergence and gradient make sense in any $\R^n$.
Curl as written is special to $\R^3$ (the cross product lives there). In $\R^2$ one
uses the \emph{scalar curl} $Q_x - P_y$, which is the $\mathbf{k}$-component above;
this is the quantity that appears in Green's theorem.

\section{The integrals}

We now assemble field $\times$ element $=$ integrand, and read off the type at each
step. In every case the final move is a parametrization that turns the integral into
an ordinary $\int(\cdots)\,\dd t$ or $\iint(\cdots)\,\dd u\,\dd v$.

\subsection{Scalar line integral: $\int_C f\,\dd s$}

\begin{signature}
$\displaystyle \int_C f\,\dd s$ \;:\; (scalar field $f$, curve $C$) $\to$ scalar.
\end{signature}

\textbf{Integrand type.} $f$ is a scalar field, evaluated along the curve it returns
a scalar; $\dd s$ is a scalar. Scalar times scalar is scalar. Integrating scalars
gives a scalar (e.g.\ the mass of a wire with density $f$).

\textbf{Unpacking to one variable.} Substitute $\dd s = \lvert\rr'(t)\rvert\,\dd t$
from Section 2.1 and evaluate $f$ at the moving point $\rr(t)$:
\[
\int_C f\,\dd s = \int_a^b f\bigl(\rr(t)\bigr)\,\lvert\rr'(t)\rvert\,\dd t .
\]
The right-hand side is an ordinary single-variable integral: integrand
$f(\rr(t))\lvert\rr'(t)\rvert$ is a plain function of $t$.

\subsection{Vector line integral (work): $\int_C \F\cdot\dd\mathbf{r}$}

\begin{signature}
$\displaystyle \int_C \F\cdot\dd\mathbf{r}$ \;:\; (vector field $\F$, oriented curve
$C$) $\to$ scalar.
\end{signature}

\textbf{Integrand type.} $\F$ is a vector field (returns a vector); $\dd\mathbf{r}$ is
a vector; the dot product collapses the two vectors to a scalar. So although both
inputs are vectors, the integrand is a \emph{scalar}, and the integral is a number
(the work done by $\F$ along $C$).

\textbf{Unpacking to one variable.} Substitute $\dd\mathbf{r}=\rr'(t)\,\dd t$:
\[
\int_C \F\cdot\dd\mathbf{r}
= \int_a^b \F\bigl(\rr(t)\bigr)\cdot\rr'(t)\,\dd t .
\]
Inside, $\F(\rr(t))\cdot\rr'(t)$ is a dot product of two concrete vectors, hence a
scalar function of $t$, and you integrate it normally.

\textbf{Connection to the scalar line integral.} Using $\dd\mathbf{r}=\TT\,\dd s$,
\[
\int_C \F\cdot\dd\mathbf{r} = \int_C (\F\cdot\TT)\,\dd s .
\]
So a work integral is just the scalar line integral of the \emph{tangential
component} $\F\cdot\TT$ (a scalar field) of the force. The vector integral and the
scalar integral are the same machine fed differently.

In components, with $\F=(P,Q,R)$ and $\dd\mathbf{r}=(\dd x,\dd y,\dd z)$, the same
integral is often written
\[
\int_C P\,\dd x + Q\,\dd y + R\,\dd z,
\]
which is literally the dot product $\F\cdot\dd\mathbf{r}$ written out term by term.

\subsection{Scalar surface integral: $\iint_S f\,\dd S$}

\begin{signature}
$\displaystyle \iint_S f\,\dd S$ \;:\; (scalar field $f$, surface $S$) $\to$ scalar.
\end{signature}

\textbf{Integrand type.} $f$ returns a scalar; $\dd S$ is a scalar; product is
scalar (e.g.\ mass of a sheet with surface density $f$).

\textbf{Unpacking to two variables.} Substitute
$\dd S = \lvert\rr_u\times\rr_v\rvert\,\dd u\,\dd v$:
\[
\iint_S f\,\dd S
= \iint_D f\bigl(\rr(u,v)\bigr)\,\lvert\rr_u\times\rr_v\rvert\,\dd u\,\dd v .
\]
The right-hand side is an ordinary double integral over the flat parameter region
$D\subset\R^2$, which you evaluate as two nested single-variable integrals.

\subsection{Flux integral: $\iint_S \F\cdot\dd\mathbf{S}$}

\begin{signature}
$\displaystyle \iint_S \F\cdot\dd\mathbf{S}$ \;:\; (vector field $\F$, oriented
surface $S$) $\to$ scalar.
\end{signature}

\textbf{Integrand type.} $\F$ returns a vector; $\dd\mathbf{S}$ is a vector; the dot
product makes the integrand a scalar (the flux: how much of $\F$ passes through $S$).

\textbf{Unpacking to two variables.} Substitute
$\dd\mathbf{S}=(\rr_u\times\rr_v)\,\dd u\,\dd v$:
\[
\iint_S \F\cdot\dd\mathbf{S}
= \iint_D \F\bigl(\rr(u,v)\bigr)\cdot(\rr_u\times\rr_v)\,\dd u\,\dd v .
\]
The integrand is a dot product of concrete vectors, a scalar function of $(u,v)$, and
the result is an ordinary double integral. Using $\dd\mathbf{S}=\nn\,\dd S$ this is
also $\iint_S(\F\cdot\nn)\,\dd S$: the scalar surface integral of the normal
component $\F\cdot\nn$.

\subsection{Volume integral: $\iiint_V f\,\dd V$}

\begin{signature}
$\displaystyle \iiint_V f\,\dd V$ \;:\; (scalar field $f$, solid $V$) $\to$ scalar.
\end{signature}
Both factors are scalars; the result is a scalar. In Cartesian coordinates it is just
three nested single-variable integrals,
$\iiint_V f\,\dd x\,\dd y\,\dd z$. There is no vector version: a region in $\R^3$ has
no leftover direction to carry, so $\dd V$ is intrinsically scalar.

\subsection{Summary of integrand types}

\begin{center}
\renewcommand{\arraystretch}{1.4}
\begin{tabular}{@{}l c c c@{}}
\toprule
\textbf{Integral} & \textbf{Field type} & \textbf{Element type} & \textbf{Integrand} \\
\midrule
$\int_C f\,\dd s$            & scalar & scalar (length) & scalar \\
$\int_C \F\cdot\dd\mathbf{r}$ & vector & vector (tangent) & scalar (via $\cdot$) \\
$\iint_S f\,\dd S$           & scalar & scalar (area)   & scalar \\
$\iint_S \F\cdot\dd\mathbf{S}$ & vector & vector (normal) & scalar (via $\cdot$) \\
$\iiint_V f\,\dd V$          & scalar & scalar (volume) & scalar \\
\bottomrule
\end{tabular}
\end{center}

Notice the pattern: every one of these integrals returns a \emph{scalar}. The vector
inputs always get reduced to a scalar integrand by a dot product before integration.
When you see ``$\cdot$'' next to a differential, read it as ``two vectors enter, one
scalar leaves.''

\section{The integral theorems: one pattern}

The four big theorems all say the same structural thing: \emph{integrating a
derivative over a region equals integrating the original field over the boundary of
that region.} Each step also drops the dimension of integration by one. This is the
multidimensional Fundamental Theorem of Calculus.

The one-dimensional FTC is the template:
\[
\int_a^b f'(x)\,\dd x = f(b) - f(a).
\]
Left side: a derivative integrated over the interval $[a,b]$. Right side: the
original function evaluated on the boundary $\{a,b\}$, which is a $0$-dimensional set
(two points, with signs). Every theorem below is this sentence with the region
upgraded.

\subsection{Gradient theorem (FTC for line integrals)}

\[
\int_C \nabla f\cdot\dd\mathbf{r} = f(\mathbf{b}) - f(\mathbf{a}),
\qquad C:\ \mathbf{a}\to\mathbf{b}.
\]
\textbf{Types.} Left: $\nabla f$ is a vector field, dotted with $\dd\mathbf{r}$ to
give a scalar line integral over the curve $C$ (a $1$-dimensional region). Right: the
scalar field $f$ evaluated on the boundary $\partial C=\{\mathbf{a},\mathbf{b}\}$
(two points, $0$-dimensional), with signs. The derivative $\nabla$ inside the
$1$-dimensional integral becomes the bare field $f$ on the $0$-dimensional boundary.

\subsection{Green's theorem (in the plane)}

\[
\oint_{\partial D} \bigl(P\,\dd x + Q\,\dd y\bigr)
= \iint_D \Bigl(\frac{\partial Q}{\partial x} - \frac{\partial P}{\partial y}\Bigr)\,\dd A .
\]
\textbf{Types.} Right: the scalar curl $Q_x-P_y$ (a scalar field) integrated over the
plane region $D$ ($2$-dimensional). Left: the vector field $(P,Q)$ integrated along
the boundary curve $\partial D$ ($1$-dimensional) as a work integral. The derivative
combination is inside on the area; the plain field is on the boundary curve.

\subsection{Stokes' theorem}

\[
\iint_S (\nabla\times\F)\cdot\dd\mathbf{S}
= \oint_{\partial S} \F\cdot\dd\mathbf{r}.
\]
\textbf{Types.} Left: $\nabla\times\F$ is a vector field (the curl), turned into a
scalar flux integrand by $\cdot\,\dd\mathbf{S}$, integrated over the surface $S$
($2$-dimensional). Right: $\F$ itself integrated along the boundary curve
$\partial S$ ($1$-dimensional) as work. Stokes is Green's theorem lifted off the
plane onto a curved surface; Green's is the special case where $S$ is a flat region
and $\F$ has no $z$-component.

\subsection{Divergence theorem}

\[
\iiint_V (\nabla\cdot\F)\,\dd V
= \iint_{\partial V} \F\cdot\dd\mathbf{S}.
\]
\textbf{Types.} Left: $\nabla\cdot\F$ is a scalar field (the divergence), integrated
over the solid $V$ ($3$-dimensional). Right: $\F$ as a flux through the boundary
surface $\partial V$ ($2$-dimensional). The derivative is inside on the volume; the
plain field flows through the boundary.

\subsection{The unifying table}

Reading down the ``inside'' column you always find a derivative of the field; reading
down the ``boundary'' column you always find the plain field. Reading across, the
dimension of the integral drops by exactly one.

\begin{center}
\renewcommand{\arraystretch}{1.5}
\begin{tabular}{@{}l c l c l@{}}
\toprule
\textbf{Theorem} & \textbf{dim} & \textbf{Inside region (derivative)} & \textbf{dim} & \textbf{On boundary (field)} \\
\midrule
FTC               & $1$ & $\displaystyle\int_a^b f'\,\dd x$ & $0$ & $f(b)-f(a)$ \\
Gradient thm      & $1$ & $\displaystyle\int_C \nabla f\cdot\dd\mathbf{r}$ & $0$ & $f(\mathbf{b})-f(\mathbf{a})$ \\
Green / Stokes    & $2$ & $\displaystyle\iint_S (\nabla\times\F)\cdot\dd\mathbf{S}$ & $1$ & $\displaystyle\oint_{\partial S}\F\cdot\dd\mathbf{r}$ \\
Divergence        & $3$ & $\displaystyle\iiint_V (\nabla\cdot\F)\,\dd V$ & $2$ & $\displaystyle\iint_{\partial V}\F\cdot\dd\mathbf{S}$ \\
\bottomrule
\end{tabular}
\end{center}

In the language of differential forms all four lines are literally one equation,
$\int_{\Omega} \dd\omega = \int_{\partial\Omega}\omega$, where $\dd$ is a single
operator (the exterior derivative) that specializes to grad, curl, and div depending
on the type of $\omega$. You do not need that machinery to use the theorems, but it
explains why the four statements rhyme: they are one statement, type-cast to
dimensions $1$, $2$, and $3$.

\section{Reading checklist}

When you meet an unfamiliar vector-calculus expression, recover its types in this
order.

\begin{enumerate}[leftmargin=1.6em]
\item \textbf{Identify the region and its dimension.} Curve ($C$, $1$D), surface
($S$, $2$D), or solid ($V$, $3$D)? The number of integral signs usually matches.
\item \textbf{Identify the differential element and its type.} $\dd s,\dd S,\dd V$
are scalars (lengths, areas, volumes). $\dd\mathbf{r},\dd\mathbf{S}$ are vectors
(they carry a direction: tangent and normal respectively).
\item \textbf{Identify the field and its type.} Lowercase $f$: scalar field.
Boldface $\F$: vector field. A ``$\cdot$'' or ``$\times$'' nearby signals a vector
field is involved.
\item \textbf{Combine and check the integrand is a scalar.} Scalar $\times$ scalar,
or vector $\cdot$ vector. If you have a vector next to a vector element with a dot,
the dot makes it scalar. The integral itself is then a number.
\item \textbf{Cast to ordinary calculus.} Parametrize: substitute
$\dd s=\lvert\rr'\rvert\,\dd t$, $\dd\mathbf{r}=\rr'\,\dd t$,
$\dd S=\lvert\rr_u\times\rr_v\rvert\,\dd u\,\dd v$, or
$\dd\mathbf{S}=(\rr_u\times\rr_v)\,\dd u\,\dd v$, evaluate the field at the moving
point, and integrate the resulting plain function of $t$ (or $u,v$).
\end{enumerate}

The single sentence to carry away: \emph{a parametrization is a type cast from a
geometric integral to the one-variable integral you already know how to do.}

\end{document}
